\section{Fork/Join --- Divisão e Conquista}
\label{sec:forkjoin}

\subsection{Implementação em Java}

\subsubsection{Descrição}

A v3 e a v4 particionam o array de pontos em exatamente $N$ fatias fixas
(uma por processador lógico) antes do início do laço de iterações. A v5
substitui esse particionamento fixo por divisão-e-conquista recursiva
com \texttt{RecursiveTask<PartialResult>}: cada tarefa recebe um
intervalo \texttt{[from, to)} do array e, se o intervalo for maior que
um limiar, divide-o ao meio, bifurca (\texttt{fork()}) a metade
esquerda para execução assíncrona, computa a metade direita
diretamente, e por fim junta (\texttt{join()}) o resultado da esquerda
para mesclar os dois parciais. Esse padrão traz \emph{work-stealing}
para o nível da tarefa: threads ociosas do \texttt{ForkJoinPool} podem
roubar subtarefas pendentes de outras threads, o que é útil quando a
carga de trabalho não está perfeitamente balanceada entre partições ---
não é o caso deste dataset sintético e uniforme, mas é um modelo
estruturalmente diferente do particionamento estático da v3/v4, e por
isso vale a pena explorar mesmo sem garantia de ganho de desempenho. O
\texttt{ScopedValue<double[][]> CENTROIDS} introduzido na v4 é mantido
sem alterações, já que se propaga corretamente para subtarefas
bifurcadas a partir de um escopo \texttt{ScopedValue.where(...).call(...)}.

O limiar de divisão é calculado como:

\begin{verbatim}
int n = Runtime.getRuntime().availableProcessors();
int threshold = points.length / (n * 4);
\end{verbatim}

ou seja, o alvo inicial é de 4 subtarefas por núcleo --- confirmado
diretamente no código-fonte da tag \texttt{v5}. Não foi realizado teste
de sensibilidade variando esse limiar: o valor é reportado como uma
escolha de design (\textit{work-stealing} só é útil se houver
subtarefas suficientes para roubar) e não como um parâmetro ajustado
empiricamente.

\begin{table}[H]
\centering
\begin{tabular}{lp{11cm}}
\hline
\textbf{Componente} & \textbf{Alteração} \\
\hline
\texttt{KMeans.java} & particionamento fixo em $N$ chunks substituído por
                       \texttt{KMeansTask extends RecursiveTask<PartialResult>}
                       (divisão recursiva por limiar, \texttt{fork()}/\texttt{join()});
                       \texttt{ForkJoinPool.commonPool()} no lugar do
                       \texttt{ExecutorService} dedicado \\
\hline
\end{tabular}
\caption{Alterações da v5 em relação à v4.}
\label{tab:v5-changes}
\end{table}

\subsubsection{Resultados}

\paragraph{Avaliação de Condição de Corrida}

Os mesmos quatro testes JCStress (JCS1--JCS4) foram executados sobre a
v5, com a mesma configuração das seções anteriores.

\begin{table}[H]
\centering
\begin{tabular}{llll}
\hline
\textbf{ID} & \textbf{Resultado} & \textbf{Outcome FORBIDDEN} & \textbf{Freq.\ máx.} \\
\hline
JCS1 & \textbf{PASSOU}            & nenhum                   & ---      \\
JCS2 & \textbf{PASSOU}            & nenhum                   & ---      \\
JCS3 & FALHOU (esperado) & \texttt{"1"} lost update em \texttt{counts[0]++}  & 28,20\% \\
JCS4 & FALHOU (esperado) & \texttt{"1.0"} lost update em \texttt{sums[0]+=1.0} & 28,96\% \\
\hline
\end{tabular}
\caption{Resultados JCStress --- regiões críticas da v5.}
\label{tab:jcstress-v5}
\end{table}

O padrão se repete: JCS1/JCS2 passam (leitura concorrente segura),
JCS3/JCS4 falham como esperado, na mesma ordem de grandeza das versões
anteriores. O modelo de divisão recursiva não introduz nenhuma nova
região crítica --- cada subtarefa folha ainda acumula em arrays locais
próprios, e a mesclagem de parciais só ocorre depois que
\texttt{join()} garante a conclusão da subtarefa filha, preservando o
mesmo padrão de confinamento local com redução sequencial das versões
anteriores.

\paragraph{Avaliação com Microbenchmark}

Os microbenchmarks foram executados com JMH 1.37 sobre JDK 25.0.1, com
\texttt{--enable-preview} (herdado da v4 para o \texttt{ScopedValue}), 1
fork, 1 iteração de aquecimento de 10 s e 3 iterações de medição de 10 s
cada.

\begin{table}[H]
\centering
\begin{tabular}{lp{4.2cm}rrrr}
\hline
\textbf{ID} & \textbf{Benchmark} & \textbf{v4} & \textbf{v5} & \textbf{Erro v5 ($\pm$)} & \textbf{Un.} \\
\hline
B1 & Execução completa (\textit{k}=10, iter=20) & 7{,}08 & 5{,}72  & 16{,}08 & s/op \\
B2 & Por iteração (\textit{k}=10, iter=5)        & 1{,}98 & 1{,}78  &  0{,}14 & s/op \\
B3 & \texttt{closestCentroid} (38 dims)           & 131{,}18 & 115{,}62 & 20{,}77 & ns/op \\
B4 & \texttt{squaredDistance} (38 dims)           &  12{,}89 &  11{,}10 &  1{,}81 & ns/op \\
\hline
\end{tabular}
\caption{Resultados JMH --- v4 vs.\ v5.}
\label{tab:jmh-v5}
\end{table}

B2 melhora levemente frente à v4 (1,98 s $\to$ 1,78 s), mas o dado mais
chamativo é a variância de B1: $\pm$16,08 s, a maior de toda a série
desde a v2-virtual ($\pm$260,68 s). O padrão é semelhante ao já
observado naquela versão --- tarefas CPU-bound que nunca bloqueiam
introduzem ruído de escalonamento quando expostas a um pool pensado
originalmente para \textit{work-stealing} de tarefas que podem ceder o
processador. Aqui o efeito é bem menor em magnitude, mas aponta na
mesma direção: a divisão recursiva com \texttt{fork()}/\texttt{join()}
tem overhead de agendamento que o particionamento estático fixo da v3/v4
não tem.

\paragraph{Avaliação com Macrobenchmark}

Os macrobenchmarks foram executados com JMeter 5.6.3 ($k=10$,
\texttt{maxIter}=3, \texttt{tol}=0,0), 120 s por configuração e ramp-up
de 30 s, com \textbf{ParallelGC} (mesma decisão de escopo da v4 ---
sem varredura de coletores para v5).

\begin{table}[H]
\centering
\begin{tabular}{rrrr}
\hline
\textbf{Threads} & \textbf{Execuções} & \textbf{Tempo (s)} & \textbf{Throughput (exec/min)} \\
\hline
1  &  93 & 120{,}2 & 46{,}4 \\
2  &  99 & 120{,}6 & 49{,}3 \\
4  & 100 & 122{,}2 & 49{,}1 \\
8  & 102 & 123{,}8 & 49{,}4 \\
16 & 112 & 128{,}3 & 52{,}4 \\
\hline
\end{tabular}
\caption{Throughput (execuções/min, ParallelGC) por número de threads JMeter --- v5.}
\label{tab:jmeter-v5}
\end{table}

O throughput da v5 é muito próximo ao da v4 em todos os níveis de carga
(faixa de 45,8--52,0 na v4 contra 46,4--52,4 na v5), ambos abaixo da v3
(52,5--62,3). Isso reforça a leitura da seção anterior: o custo adicional
de coordenação --- seja via \texttt{StructuredTaskScope} (v4) ou via
\texttt{ForkJoinPool} com divisão recursiva (v5) --- supera, neste
dataset uniforme e de carga bem distribuída, qualquer ganho estrutural
que \textit{work-stealing} poderia oferecer.

\paragraph{Avaliação de Profile}

O profile foi capturado com JFR (\texttt{settings=profile}) sobre uma
execução completa do algoritmo. A gravação totalizou 2.452 amostras de
\texttt{jdk.ExecutionSample}.

\begin{table}[H]
\centering
\small
\begin{tabular}{p{7.4cm}rrp{3cm}}
\hline
\textbf{Método} & \textbf{Amostras} & \textbf{CPU (\%)} & \textbf{Categoria} \\
\hline
\texttt{KMeans.squaredDistance}       & 1084 & 44{,}2 & Algoritmo \\
\texttt{KMeans.closestCentroid}       &  253 & 10{,}3 & Algoritmo \\
\texttt{FloatingDecimal.copyDigits}   &  193 &  7{,}9 & Parsing CSV (inicialização) \\
\texttt{String.split}                 &  185 &  7{,}5 & Parsing CSV (inicialização) \\
\texttt{FloatingDecimal.parseDouble}  &  177 &  7{,}2 & Parsing CSV (inicialização) \\
\texttt{FloatingDecimal.readJavaFormatString} &  138 &  5{,}6 & Parsing CSV (inicialização) \\
\texttt{FloatingDecimal.doubleValue}  &   78 &  3{,}2 & Parsing CSV (inicialização) \\
\texttt{BufferedReader.readLine}      &   51 &  2{,}1 & I/O \\
\hline
\end{tabular}
\caption{Hot methods na v5 (JFR, 2.452 amostras).}
\label{tab:jfr-hot-v5}
\end{table}

O algoritmo (\texttt{squaredDistance} + \texttt{closestCentroid}) domina
com \textasciitilde55\% das amostras --- proporção maior que na v4
(\textasciitilde46\%), sinal de que a divisão recursiva tem menos
overhead de coordenação visível no profile que o
\texttt{StructuredTaskScope}. O \textit{parsing} de CSV da leitura única
inicial (\texttt{loadAllPoints}) segue presente pelo mesmo motivo já
discutido na v4: a janela de gravação do JFR cobre a inicialização,
mesmo que ela seja excluída da medição JMH/JMeter. O arquivo
\texttt{io.txt} para a v5 não registra nenhum evento
\texttt{jdk.FileRead}, diferente da v4 (que tinha 17 linhas de eventos
de \textit{classloading}, sem leitura de dataset em nenhuma das duas).

Foram registradas 56 coletas de GC, com pausa total de 978,19 ms e
pausa média de 17,47 ms por evento. O heap atingiu pico de 2.560,0 MB
(\textasciitilde2,56 GB), próximo dos \textasciitilde2,7 GB e
\textasciitilde2,76 GB de v3 e v4 --- esperado, pelo mesmo dataset em
memória.

\subsection{Discussão}

A Tabela~\ref{tab:jmh-comparativo-v5} estende o comparativo com a v5.

\begin{table}[H]
\centering
\begin{tabular}{lrrrr}
\hline
\textbf{Abordagem} & \textbf{B1 (s/op)} & \textbf{B2 (s/op)} & \textbf{Erro B1} & \textbf{Erro B2} \\
\hline
Serial & 212{,}57 & 60{,}05 & $\pm$34{,}48 & $\pm$2{,}46 \\
v3     &   4{,}55 &  1{,}60 & $\pm$17{,}98 & $\pm$0{,}17 \\
v4     &   7{,}08 &  1{,}98 &  $\pm$1{,}26 & $\pm$0{,}07 \\
v5     &   5{,}72 &  1{,}78 & $\pm$16{,}08 & $\pm$0{,}14 \\
\hline
\end{tabular}
\caption{Comparativo JMH: Serial, v3, v4 e v5 (s/op).}
\label{tab:jmh-comparativo-v5}
\end{table}

A v5 é, como a v4, uma variação estrutural na camada de coordenação, não
uma otimização de desempenho: os números de B2 e de throughput ficam
muito próximos aos da v4, ambos abaixo da v3. A diferença pedagógica
está no modelo --- divisão-e-conquista recursiva com \textit{work-stealing}
em vez de particionamento estático --- que só demonstraria vantagem real
em cargas desbalanceadas entre partições, o que não é o caso deste
dataset sintético e uniformemente distribuído. O JCStress confirma, mais
uma vez, que o confinamento local por tarefa permanece correto
independentemente de como as tarefas são criadas e distribuídas: JCS1/JCS2
passam, JCS3/JCS4 falham exatamente como esperado.

A variância elevada de B1 ($\pm$16,08 s) é o achado mais interessante
desta seção: mesmo em uma escala bem menor que a da v2-virtual
($\pm$260,68 s), ela reforça a mesma lição --- pools de tarefas
pensados para \textit{work-stealing} introduzem ruído de escalonamento
em cargas CPU-bound que nunca bloqueiam, e esse efeito aparece tanto em
virtual threads sobre \texttt{ForkJoinPool} (v2-virtual) quanto em
divisão recursiva explícita sobre o mesmo framework (v5).
